\textbf{Phần 1. Trắc nghiệm lựa chọn}
\setcounterpageref{ex}{0}

\begin{ex} %[3H1V20]
Để vẽ được một hình tròn có bán kính $4\text{ cm}$ trên giấy, ta cần chỉnh khoảng cách giữa đầu kim và đầu bút chì của compa bằng
\choice
{$8\text{ cm}$}
{$2\text{ cm}$}
{$16\text{ cm}$}
{\True $4\text{ cm}$}
\loigiai{
Khoảng cách giữa hai đầu của compa đúng bằng bán kính hình tròn cần vẽ, tức là $4\text{ cm}$.
\textbf{Đáp án của câu hỏi là D}
}
\end{ex}

\begin{ex} %[3H1V20]
Bạn Nam dùng các que tính dài bằng nhau để ghép thành 3 hình tam giác riêng biệt không chung cạnh. Nam cần dùng ít nhất bao nhiêu que tính?
\choice
{$6$ que tính}
{\True $9$ que tính}
{$7$ que tính}
{$8$ que tính}
\loigiai{
Mỗi hình tam giác cần 3 que tính, 3 tam giác riêng biệt cần: $3 \times 3 = 9$ que tính.
\textbf{Đáp án của câu hỏi là B}
}
\end{ex}

\begin{ex} %[3H1V20]
Nếu cho phép các hình tam giác có cạnh chung với nhau, hỏi dùng ít nhất bao nhiêu que tính dài bằng nhau để ghép được 2 hình tam giác?
\choice
{$6$ que tính}
{$4$ que tính}
{\True $5$ que tính}
{$3$ que tính}
\loigiai{
Hai hình tam giác ghép chung 1 cạnh ở giữa sẽ dùng: $3 + 3 - 1 = 5$ que tính.
\textbf{Đáp án của câu hỏi là C}
}
\end{ex}

\begin{ex} %[3H1V20]
Để vẽ một góc vuông đỉnh $O$ bằng ê-ke, bước nào sau đây là chuẩn xác nhất?
\choice
{\True Đặt đỉnh góc vuông của ê-ke trùng với điểm $O$, vẽ hai cạnh theo hai mép góc vuông của thước}
{Đặt mép thước bất kì vào điểm $O$}
{Vẽ một đường tròn rồi lấy 1 phần 4}
{Dùng thước thẳng vẽ ngẫu nhiên hai đoạn thẳng}
\loigiai{
Cách vẽ góc vuông chuẩn: đặt đỉnh góc vuông ê-ke vào đỉnh góc, vạch 2 tia theo 2 cạnh góc vuông của ê-ke.
\textbf{Đáp án của câu hỏi là A}
}
\end{ex}

\begin{ex} %[3H1V20]
Người ta gấp một mảnh giấy hình vuông làm đôi theo đường chéo, sau đó gấp đôi tiếp một lần nữa theo đường cao. Mảnh giấy nhận được có hình dạng là
\choice
{Hình vuông nhỏ}
{Hình chữ nhật}
{Hình tứ giác không đều}
{\True Hình tam giác có 1 góc vuông}
\loigiai{
Gấp đôi hình vuông theo đường chéo được tam giác vuông cân. Gấp đôi lần nữa vẫn được hình tam giác vuông cân nhỏ hơn.
\textbf{Đáp án của câu hỏi là D}
}
\end{ex}

\begin{ex} %[3H1V20]
Một hoa văn trang trí hình vuông được chia thành 4 ô vuông nhỏ bằng nhau. Người ta tô màu đỏ vào 2 ô đối đỉnh, 2 ô còn lại tô màu xanh. Phần ô vuông màu đỏ chiếm bao nhiêu phần của hoa văn?
\choice
{Một phần tư ($\dfrac{1}{4}$)}
{\True Một phần hai ($\dfrac{1}{2}$)}
{Ba phần tư ($\dfrac{3}{4}$)}
{Một phần ba ($\dfrac{1}{3}$)}
\loigiai{
Có tất cả 4 ô vuông nhỏ, 2 ô màu đỏ. Vậy phần màu đỏ chiếm $2 : 4 = \dfrac{1}{2}$ hình vuông hoa văn.
\textbf{Đáp án của câu hỏi là B}
}
\end{ex}

\begin{ex} %[3H1V20]
Dùng 12 que diêm bằng nhau để xếp thành một hình chữ nhật có chu vi lớn nhất có thể. Hình chữ nhật đó có chiều dài gồm mấy que diêm?
\choice
{$4$ que diêm}
{$3$ que diêm}
{\True $5$ que diêm}
{$6$ que diêm}
\loigiai{
Nửa chu vi gồm $12 : 2 = 6$ que diêm. Chiều dài lớn nhất có thể là $6 - 1 = 5$ que diêm (khi chiều rộng là 1 que diêm).
\textbf{Đáp án của câu hỏi là C}
}
\end{ex}

\begin{ex} %[3H1V20]
Trong một hình vẽ gồm một đường tròn tâm $O$ và một hình vuông nằm bên trong sao cho 4 đỉnh của hình vuông đều nằm trên đường tròn. Khi đó tâm $O$ của hình tròn nằm ở vị trí nào của hình vuông?
\choice
{\True Giao điểm của hai đường chéo hình vuông}
{Tại một đỉnh bất kì của hình vuông}
{Tại trung điểm của một cạnh hình vuông}
{Nằm ngoài hình vuông}
\loigiai{
Hình tròn đi qua 4 đỉnh của hình vuông có tâm chính là giao điểm của hai đường chéo hình vuông.
\textbf{Đáp án của câu hỏi là A}
}
\end{ex}

\begin{ex} %[3H1V20]
Muốn vẽ 4 hình tròn bằng nhau có các tâm lần lượt là 4 đỉnh của một hình chữ nhật sao cho mỗi hình tròn tiếp xúc với 2 hình tròn liền kề. Điều kiện của hình chữ nhật đó là
\choice
{Hình chữ nhật có chiều dài gấp đôi chiều rộng}
{Hình chữ nhật có chiều dài gấp 3 chiều rộng}
{Bất kì hình chữ nhật nào cũng vẽ được}
{\True Hình chữ nhật đó phải là hình vuông}
\loigiai{
Để 4 hình tròn bằng nhau tiếp xúc đôi một theo cạnh thì khoảng cách giữa các đỉnh liên tiếp phải bằng nhau (bằng 2 lần bán kính), do đó các cạnh phải bằng nhau, tức là hình vuông.
\textbf{Đáp án của câu hỏi là D}
}
\end{ex}

\begin{ex} %[3H1V20]
Bạn Hùng cắt một tờ giấy hình tròn thành 4 phần bằng nhau bằng 2 nhát cắt thẳng đi qua tâm. Góc ở tâm của mỗi phần giấy nhận được là
\choice
{Góc nhọn}
{\True Góc vuông}
{Góc tù}
{Góc bẹt}
\loigiai{
Hai nhát cắt thẳng qua tâm chia hình tròn thành 4 phần bằng nhau thì hai đường cắt vuông góc với nhau, góc ở mỗi phần là góc vuông ($360^\circ : 4 = 90^\circ$).
\textbf{Đáp án của câu hỏi là B}
}
\end{ex}

\begin{ex} %[3H1V20]
Một họa sĩ vẽ một bức tranh ghép từ các hình tam giác đều có cạnh $2\text{ cm}$. Hàng thứ nhất có 1 tam giác, hàng thứ hai có 3 tam giác, hàng thứ ba có 5 tam giác. Hỏi bức tranh gồm tất cả bao nhiêu hình tam giác đều nhỏ?
\choice
{$8$ hình tam giác}
{$10$ hình tam giác}
{\True $9$ hình tam giác}
{$12$ hình tam giác}
\loigiai{
Tổng số hình tam giác nhỏ là: $1 + 3 + 5 = 9$ hình tam giác.
\textbf{Đáp án của câu hỏi là C}
}
\end{ex}

\begin{ex} %[3H1V20]
Dùng thước và compa vẽ một đường tròn tâm $O$ bán kính $5\text{ cm}$. Trên đường tròn lấy điểm $A$. Khoảng cách từ điểm $A$ đến tâm $O$ đo được là
\choice
{\True $5\text{ cm}$}
{$10\text{ cm}$}
{$2\text{ cm và } 5\text{ mm}$}
{$0\text{ cm}$}
\loigiai{
Mọi điểm nằm trên đường tròn đều cách tâm một khoảng bằng bán kính, tức là $5\text{ cm}$.
\textbf{Đáp án của câu hỏi là A}
}
\end{ex}

\textbf{Phần 2. Trắc nghiệm đúng sai}
\setcounterpageref{ex}{0}

\begin{ex} %[3H1V20]
Xét tính đúng sai của các phát biểu về vẽ hình học phẳng:
\choiceTF
{\True Khi dùng compa vẽ đường tròn, chân trụ có kim phải được giữ cố định tại tâm.}
{Chỉ dùng một thước thẳng có vạch đo có thể vẽ chính xác góc vuông mà không cần ê-ke.}
{\True Dùng 4 que tính bằng nhau có thể xếp thành một hình vuông hoặc một hình thoi.}
{Mọi hình vẽ trang trí đều phải có góc vuông.}
\loigiai{
\begin{itemize}
    \item Ý a) Đúng: Quy tắc dùng compa.
    \item Ý b) Sai: Sai vì thước thẳng không đo góc trực tiếp được góc vuông.
    \item Ý c) Đúng: Đúng vì 4 cạnh bằng nhau.
    \item Ý d) Sai: Sai vì có nhiều hoa văn tròn, uốn lượn không có góc vuông.
\end{itemize}
\begin{itemize}
    \item a) \textbf{Đúng}.
    \item b) \textbf{Sai}.
    \item c) \textbf{Đúng}.
    \item d) \textbf{Sai}.
\end{itemize}
}
\end{ex}

\begin{ex} %[3H1V20]
Một tờ giấy hình chữ nhật kích thước $12\text{ cm} \times 8\text{ cm}$ được cắt ra để trang trí:
\choiceTF
{\True Hình vuông lớn nhất có thể cắt ra từ tờ giấy có cạnh bằng $8\text{ cm}$.}
{\True Phần còn lại sau khi cắt hình vuông lớn nhất là một hình chữ nhật kích thước $8\text{ cm} \times 4\text{ cm}$.}
{Có thể cắt tờ giấy thành 3 hình vuông bằng nhau mà không thừa mảnh nào.}
{\True Từ phần còn lại ($8\text{ cm} \times 4\text{ cm}$) ta cắt tiếp được đúng 2 hình vuông cạnh $4\text{ cm}$.}
\loigiai{
\begin{itemize}
    \item Ý a) Đúng: Cạnh hình vuông lớn nhất bằng chiều rộng.
    \item Ý b) Đúng: $12 - 8 = 4\text{ cm}$.
    \item Ý c) Sai: Sai vì diện tích 96 không chia hết cho 3 hình vuông nguyên cạnh.
    \item Ý d) Đúng: $8 : 4 = 2$ hình vuông.
\end{itemize}
\begin{itemize}
    \item a) \textbf{Đúng}.
    \item b) \textbf{Đúng}.
    \item c) \textbf{Sai}.
    \item d) \textbf{Đúng}.
\end{itemize}
}
\end{ex}

\textbf{Phần 3. Trả lời ngắn}
\setcounterpageref{ex}{0}

\begin{ex} %[3H1V20]
Cần ít nhất bao nhiêu que diêm bằng nhau để ghép thành 3 hình vuông riêng biệt không chung cạnh?
\shortans{12}
\loigiai{
\textbf{Đáp số:} $12$\\
Mỗi hình vuông cần 4 que diêm: $3 \times 4 = 12$ que.
}
\end{ex}

\begin{ex} %[3H1V20]
Dùng các que tính bằng nhau để xếp thành 2 hình vuông có chung 1 cạnh. Hỏi cần dùng tất cả bao nhiêu que tính?
\shortans{7}
\loigiai{
\textbf{Đáp số:} $7$\\
Hai hình vuông ghép chung 1 cạnh cần: $4 + 4 - 1 = 7$ que tính.
}
\end{ex}

\begin{ex} %[3H1V20]
Để vẽ một hình tròn có đường kính $16\text{ cm}$, bạn Mai cần mở rộng hai đầu compa một khoảng cách bằng bao nhiêu xăng-ti-mét?
\shortans{8}
\loigiai{
\textbf{Đáp số:} $8$\\
Bán kính compa cần mở là: $16 : 2 = 8\text{ (cm)}$.
}
\end{ex}

\begin{ex} %[3H1V20]
Một đường gấp khúc gồm 5 đoạn thẳng bằng nhau có độ dài mỗi đoạn là $6\text{ cm}$. Tính tổng độ dài đường gấp khúc đó theo xăng-ti-mét.
\shortans{30}
\loigiai{
\textbf{Đáp số:} $30$\\
Độ dài đường gấp khúc: $5 \times 6 = 30\text{ (cm)}$.
}
\end{ex}

\begin{ex} %[3H1V20]
Một dải băng giấy dài $36\text{ cm}$ được gấp làm 4 đoạn bằng nhau dích dắc để tạo hình trang trí. Hỏi mỗi đoạn gấp dài bao nhiêu xăng-ti-mét?
\shortans{9}
\loigiai{
\textbf{Đáp số:} $9$\\
Độ dài mỗi đoạn gấp là: $36 : 4 = 9\text{ (cm)}$.
}
\end{ex}

\begin{ex} %[3H1V20]
Dùng 6 hình tam giác đều nhỏ bằng nhau ghép lại có thể tạo thành một hình lục giác đều. Nếu mỗi tam giác nhỏ có chu vi $18\text{ cm}$ thì độ dài cạnh của lục giác ghép được là bao nhiêu xăng-ti-mét?
\shortans{6}
\loigiai{
\textbf{Đáp số:} $6$\\
Cạnh của tam giác đều nhỏ là: $18 : 3 = 6\text{ cm}$. Cạnh của lục giác đều chính bằng cạnh của tam giác đều nhỏ, bằng $6\text{ cm}$.
}
\end{ex}
